% Section 4: Method (about 1.5 pages). Plan: PAPER_PLAN 4 / §4.
%   4.1 The check-in-level representation (Check2HGI): four-level graph (check-in, place, region, city);
%       infomax objective; one vector per check-in. Fig. \ref{fig:dataflow} (top priority).
%   4.2 The single model: cross-attention two-stream stack + private spatial path; one model, one forward
%       pass, two predictions. Fixed-weight sum (0.50/0.50); region output plain cross-entropy,
%       category output cross-entropy with logit adjustment (tau=0.5, train time only), which
%       replaces class weighting. The dedicated cat model carries the same adjustment.
%       ⚠ COST (2026-07-08 audit): the joint model is LARGER than the two Table-3 dedicated models COMBINED
%       (~2.6-4.2x params, ~4x ops, measured on the instantiated v17 model). The old "+5 percent" and the
%       INV2 "17-20% fewer params / 28% fewer ops" numbers used a different basis (single-task copies of the
%       joint architecture itself) and must NEVER return to the prose. The honest claim is operational
%       (one artifact, one forward pass), not arithmetic. Fig. \ref{fig:model}.
\section{Method}
\label{sec:method}

\subsection{The check-in-level representation}
\label{sec:method-rep}

We want each visit, not each place, to have its own vector; we call the resulting representation Check2HGI (Section~\ref{sec:related-embeddings}).
Figure~\ref{fig:dataflow} shows the full data flow, from check-ins to the two outputs.

We build a graph with four levels: the check-in, its place, the place's region (a census tract), and the city. Edges connect each level to the one above it and link a user's consecutive check-ins with a weight that decays as the time gap between the visits grows; same-place visits connect through their shared place node one level up, and nearby places are linked at the place level. Each visit's category, time of day, and day of week enter as input features of its node, not as edges. Four elapsed-time features join them: the time since the user's previous visit and since the user's first visit, both on a logarithmic scale, the gap to the previous visit within the same day, and an indicator for a user's first visit. Each is measured up to the visit itself, so a node describes the visit and the history preceding it. These give the representation a sense of tempo, distinguishing a visit that follows minutes after the last one from a visit that opens a new outing, which a categorical hour-and-weekday encoding alone cannot express. This keeps a hierarchical graph's place-to-region-to-city geography and adds the check-in as a bottom level. We train the graph mainly with an infomax objective, under which each vector learns to match its real neighborhood and reject a shuffled one~\cite{velickovic2019dgi,huang2023hgi}. Two small label-free auxiliary terms are added (weights 0.3 and 0.1): a masked reconstruction of each place's aggregated category
features, and an anchor to a place embedding pre-trained, label-free,
on the same data. The training
uses no task label: it
never sees the next category or the next region. From the trained graph, we extract one 64-dimensional vector per check-in (its region nodes have trained vectors as well), so a model sees a sequence of per-visit vectors rather than repeated per-place ones.

\subsection{The single multi-task model}
\label{sec:method-model}

We then train one model that reads a window of recent per-visit vectors and predicts the next visit's category and region at once. Figure~\ref{fig:model} shows the design. Each task has its own input. The category task reads the window of per-visit vectors (the semantic stream); the region task reads the same window of visits, each visit now represented by the trained vector of its region node from the same graph (the spatial stream). Both pass through private
per-task encoders (a small input network per task, with no shared weights) into the shared trunk, a cross-attention stack of two blocks, so each prediction still uses both inputs. In each block, attention lets each
stream read the other's features, while each keeps its own feed-forward
weights. The tasks therefore share by exchanging information between per-task
streams, not by owning hidden layers in common. The trunk then feeds a category output and a region output; the region output also keeps a private spatial path: a small branch inside the one model (not a second model) that reads the spatial input window, bypassing the shared trunk. The category task does not touch this branch. Together, the region task's encoder, its output, and this private path form the region pathway.

The training loss is a fixed-weight sum of the two outputs,
\begin{equation}\label{eq:loss}
L = 0.5\,L_{\mathrm{cat}} + 0.5\,L_{\mathrm{reg}},
\end{equation}
where $L_{\mathrm{reg}}$ is a plain unweighted cross-entropy loss and $L_{\mathrm{cat}}$ is a
cross-entropy loss with logit adjustment: $\tau\log P_{\mathrm{train}}(y)$, with $\tau=0.5$, is added
to the category logits during training only, which shifts the decision boundary toward the balanced
posterior that macro-F1 rewards. Inference uses the unadjusted logits. The weights and $\tau$ were
selected once on validation during development and held fixed across all six datasets. Class
weighting, tested on both outputs, lowered both region accuracy and category macro-F1, and logit
adjustment replaces it on the category output. This is standard fixed-weight joint
training~\cite{caruana1997multitask}, kept simple by design so that any improvement over the
dedicated single-task models comes from the shared representation, not from an adaptive weighting
scheme. The dedicated category model receives the same adjustment, so the comparison is not affected
by it.

The joint model includes the shared trunk and both task outputs, so it is
larger than either dedicated model, and a forward pass costs more compute than running the
two small dedicated models: the joint model has about 4.2 million parameters at
Alabama against 1.1 million for the two dedicated models combined (5.2 against
2.0 at California). What the single model provides is operational rather
than arithmetic: one artifact to train, version, and deploy, and one forward
pass whose one set of inputs produces both answers at once.
